\documentclass[11pt, a4paper]{article}

% ---------- Packages ----------
\usepackage[utf8]{inputenc}
\usepackage[T1]{fontenc}
\usepackage[margin=1in]{geometry}
\usepackage{graphicx}
\usepackage{amsmath, amssymb}
\usepackage{booktabs}
\usepackage{caption}
\usepackage{subcaption}
\usepackage{hyperref}
\usepackage[round]{natbib}
\usepackage{siunitx}
\usepackage{float}
\usepackage{fancyhdr}

\hypersetup{
    colorlinks=true,
    linkcolor=black,
    citecolor=blue,
    urlcolor=blue
}

\pagestyle{fancy}
\fancyhf{}
\fancyhead[L]{\nouppercase{\leftmark}}
\fancyfoot[C]{\thepage}

% ---------- Title information ----------
\title{Effects of Short-Term Sleep Deprivation on Working Memory Recall}
\author{
    Jordan Reyes\thanks{Corresponding author: \texttt{j.reyes@example.edu}} \\
    Department of Cognitive Science, Example University
    \and
    Priya Nandakumar \\
    Department of Cognitive Science, Example University
}
\date{\today}

\begin{document}

\maketitle

\begin{abstract}
Sleep loss is widely reported to impair cognitive performance, but the
magnitude of its effect on working memory specifically remains debated. We
tested 32 healthy adults on a digit-span recall task under two conditions:
a normal-sleep night (\qty{7}{}--\qty{9}{} hours) and a restricted-sleep
night (\qty{4}{} hours), using a within-subjects, counterbalanced design.
Participants in the restricted-sleep condition recalled significantly
fewer digits on average ($M = 5.8$, $SD = 1.1$) than in the normal-sleep
condition ($M = 7.2$, $SD = 1.3$), $t(31) = 6.42$, $p < 0.001$, Cohen's
$d = 1.13$. These results support the hypothesis that even a single night
of restricted sleep produces a measurable, moderate-to-large impairment in
working memory recall. We discuss implications for shift workers and
students, and note limitations related to sample size and self-reported
sleep compliance.
\end{abstract}

\noindent\textbf{Keywords:} sleep deprivation, working memory, cognitive
performance, digit span, within-subjects design

\tableofcontents
\newpage

% ---------- Introduction ----------
\section{Introduction}
\label{sec:introduction}

Working memory -- the capacity to hold and manipulate information over
short intervals -- underlies everyday tasks ranging from mental arithmetic
to following multi-step instructions. Prior research has linked sleep loss
to broad declines in attention and reaction time \citep{killgore2010}, but
findings on working memory specifically have been mixed, with effect sizes
ranging from negligible to large depending on task design and the degree
of sleep restriction \citep{lowe2017}.

This report addresses a narrower question: does a single night of
moderate sleep restriction (4 hours) measurably impair performance on a
standard digit-span recall task, relative to a normal night of sleep, in
healthy young adults? We hypothesized that restricted sleep would reduce
mean digit-span recall relative to the normal-sleep condition.

The remainder of this report is organized as follows.
Section~\ref{sec:related-work} reviews related work.
Section~\ref{sec:methods} describes the experimental design and analysis.
Section~\ref{sec:results} presents the results.
Section~\ref{sec:discussion} interprets the findings and their
limitations, and Section~\ref{sec:conclusion} concludes.

% ---------- Related Work ----------
\section{Related Work}
\label{sec:related-work}

\citet{killgore2010} reviewed evidence that sleep deprivation degrades
attention, reaction time, and higher-order decision making, but noted that
working-memory-specific effects were less consistently reported across
studies, in part due to differing task demands. \citet{lowe2017} conducted
a meta-analysis of 60 studies and found a small-to-moderate pooled effect
of total sleep deprivation on working memory tasks ($g \approx 0.35$),
with larger effects for partial sleep restriction protocols closer to the
one used here. Our study extends this literature with a controlled,
within-subjects comparison using a single restriction night rather than
total deprivation, which is more representative of common real-world sleep
loss (e.g., exam periods, shift work transitions).

% ---------- Methods ----------
\section{Methods}
\label{sec:methods}

\subsection{Participants}
Thirty-two healthy adults (18 female, 14 male; age $M = 22.4$,
$SD = 2.7$) with no reported sleep disorders were recruited from the
university community. All participants provided informed consent, and the
protocol was approved by the university's institutional review board.

\subsection{Experimental Setup}
A within-subjects, counterbalanced design was used. Each participant
completed the digit-span recall task twice, one week apart: once after a
normal-sleep night (self-reported \qty{7}{}--\qty{9}{} hours, verified by
wrist actigraphy) and once after a restricted-sleep night (sleep window
capped at \qty{4}{} hours, also verified by actigraphy). Condition order
was counterbalanced across participants. The task itself was administered
each morning between 9:00 and 10:00 to control for time-of-day effects.

\subsection{Data Analysis}
The primary outcome was the longest correctly recalled digit sequence,
averaged across 10 trials per session. We modeled the difference between
conditions with a paired-samples $t$-test. The within-subject effect can
be summarized as
\begin{equation}
    \label{eq:diff}
    d_i = y_i^{\text{normal}} - y_i^{\text{restricted}}
\end{equation}
where $y_i^{\text{normal}}$ and $y_i^{\text{restricted}}$ are participant
$i$'s mean recall scores in the normal- and restricted-sleep conditions,
respectively, and $d_i$ is the within-subject difference used in the
paired $t$-test.

% ---------- Results ----------
\section{Results}
\label{sec:results}

Mean digit-span recall was lower in the restricted-sleep condition
($M = 5.8$, $SD = 1.1$) than in the normal-sleep condition ($M = 7.2$,
$SD = 1.3$). The paired-samples $t$-test confirmed this difference was
statistically significant, $t(31) = 6.42$, $p < 0.001$, with a large
effect size (Cohen's $d = 1.13$). Table~\ref{tab:results} summarizes the
descriptive statistics, and Figure~\ref{fig:recall} shows the
distribution of individual differences.

\begin{figure}[H]
    \centering
    % \includegraphics[width=0.7\textwidth]{figures/recall-difference.pdf}
    \caption{Distribution of within-subject differences in digit-span
    recall (normal sleep minus restricted sleep) across 32 participants.
    Positive values indicate better recall under normal sleep. The
    distribution is shifted right of zero, consistent with the observed
    group-level effect.}
    \label{fig:recall}
\end{figure}

\begin{table}[H]
    \centering
    \caption{Descriptive statistics for digit-span recall by condition.}
    \label{tab:results}
    \begin{tabular}{lccc}
        \toprule
        Condition        & Mean & SD  & $n$ \\
        \midrule
        Normal sleep     & 7.2  & 1.3 & 32  \\
        Restricted sleep & 5.8  & 1.1 & 32  \\
        \bottomrule
    \end{tabular}
\end{table}

No significant effect of condition order was found ($p = 0.41$),
suggesting the counterbalancing successfully controlled for practice
effects.

% ---------- Discussion ----------
\section{Discussion}
\label{sec:discussion}

The observed drop of roughly 1.4 digits between conditions is consistent
with the moderate-to-large effects reported for partial sleep restriction
in \citet{lowe2017}, and supports the hypothesis that even a single night
of reduced sleep meaningfully impairs working memory recall. Practically,
a difference of this size could affect performance on tasks that place
heavy demands on short-term retention, such as mental arithmetic,
multi-step procedures, or exam recall.

Several limitations should be noted. The sample was drawn from a single
university population, limiting generalizability to other age groups or
occupational settings (e.g., shift workers with chronic partial sleep
restriction). Sleep was restricted under laboratory instructions rather
than occurring naturally, which may not fully capture the psychological
and physiological profile of real-world sleep loss. Finally, only one
cognitive measure (digit span) was used; broader working-memory batteries
would strengthen the generality of the conclusion.

% ---------- Conclusion ----------
\section{Conclusion}
\label{sec:conclusion}

A single night of sleep restricted to 4 hours produced a statistically
significant and practically meaningful reduction in working memory recall
relative to a normal night of sleep, in a sample of healthy young adults.
Future work should extend this design to repeated nights of partial sleep
restriction and to broader working-memory tasks to assess how the effect
accumulates and generalizes.

% ---------- Acknowledgments ----------
\section*{Acknowledgments}
We thank the volunteers who participated in this study and the university
sleep laboratory staff for their assistance with actigraphy monitoring.

% ---------- Bibliography ----------
\bibliographystyle{plainnat}
\bibliography{model_references}

% ---------- Appendix ----------
\appendix
\section{Supplementary Material}
\label{sec:appendix}

Raw per-participant recall scores and actigraphy logs are available from
the corresponding author upon request.

\end{document}
